\documentclass[preprint,12pt]{elsarticle}
%% Packages
\usepackage{float}
\usepackage{amssymb}
\usepackage{amsmath}
\usepackage{setspace}
\usepackage{graphicx}
\usepackage{booktabs}
\usepackage{enumitem}
\usepackage{algorithm}
\usepackage{subcaption}
\usepackage{algpseudocode}
\usepackage[table]{xcolor}
\usepackage[colorlinks=true,citecolor=blue,linkcolor=black]{hyperref}
\usepackage[textheight=9in,textwidth=6.5in]{geometry}

\setstretch{0.98}

%% Lists
\setlist[itemize]{
    topsep=2pt,
    itemsep=2pt,
    parsep=0pt,
    partopsep=0pt
}

%% Paragraph formatting
\setlength{\parindent}{0pt}
\setlength{\parskip}{0.4em}


%% Custom commands for titles
\makeatletter

\renewcommand\section{\@startsection
  {section}{1}{\z@}
  {-3.5ex \@plus -1ex \@minus -.2ex}
  {2.3ex \@plus.2ex}
  {\normalfont\Large\bfseries}}

\renewcommand\subsection{\@startsection
  {subsection}{2}{\z@}
  {-3.25ex\@plus -1ex \@minus -.2ex}
  {1.5ex \@plus .2ex}
  {\normalfont\large\bfseries}}

\renewcommand\subsubsection{\@startsection
  {subsubsection}{3}{\z@}
  {-2.5ex\@plus -1ex \@minus -.2ex}
  {1ex \@plus .2ex}
  {\normalfont\normalsize\bfseries}}

\makeatother

%% ==================================================================== 
%% Journal
%% ====================================================================
\journal{Computers \& Industrial Engineering}
\biboptions{sort&compress}

%% ==================================================================== 
%% Document 
%% ====================================================================
\begin{document}
\begin{frontmatter}

\title{ A Multi-Level Learning Framework for Adaptive Search Control in the VRPTW-S }

%% ==================================================================== 
%% Authors
%% ====================================================================
\author[1]{Nada Farah Djedjig}
\ead{nadadjedjig@gmail.com}

\author[1]{Supervisor Name}
\ead{supervisor@gmail.com}

\author[1]{Supervisor Name}
\ead{supervisor@gmail.com}

\address[1]{Department of Computer Science, University Name, City, Country}

%% ==================================================================== 
%% Abstract
%% ==================================================================== 
\begin{abstract}
\noindent
The Vehicle Routing Problem with Time Windows and Synchronization constraints (VRPTW-S) is a challenging variant of the classical vehicle routing problem arising in home healthcare and collaborative service operations. In addition to routing and scheduling decisions, the problem requires the coordinated execution of synchronized visits, creating strong temporal dependencies between routes and significantly increasing computational complexity. Although recent metaheuristics, particularly Generalized Variable Neighborhood Search (GVNS), have achieved competitive results on benchmark instances, their performance still relies on static search policies, including fixed neighborhood orders, predefined local search strategies, and offline parameter calibration.

This paper proposes a learning-guided adaptive GVNS framework for the VRPTW-S, introducing adaptive control mechanisms at multiple levels of the search process. First, reinforcement learning is used to dynamically select neighborhood operators during the Variable Neighborhood Descent phase through both tabular Q-learning and Deep Q-Network (DQN) approaches. Second, a two-head Random Forest model is introduced to adaptively choose between First Improvement and Best Improvement local search strategies according to the current search context. Third, a Proximal Policy Optimization (PPO) controller is used to dynamically regulate key search hyperparameters online.

Experiments conducted on Solomon-derived synchronized benchmark instances ranging from 50 to 400 customers demonstrate that learning-guided adaptation improves both scalability and solution quality compared with static GVNS configurations. In particular, DQN-based operator selection exhibits superior robustness on large-scale instances, while hybrid adaptive strategies provide additional gains by improving search diversification and intensification balance.
\end{abstract}


%% ==================================================================== 
%% Keywords
%% ==================================================================== 
\begin{keyword}
\small
Vehicle Routing Problem \sep 
Synchronization Constraints \sep 
Variable Neighborhood Search \sep 
Reinforcement Learning \sep 
Deep Q-Network \sep 
Proximal Policy Optimization \sep 
Adaptive Metaheuristics
\end{keyword}

\end{frontmatter}
% \newpage

%% ==================================================================== 
%% 1. INTRODUCTION 
%% ====================================================================
\section{Introduction}
\label{sec:intro}

%% ==================================================================== 
%% 2. Related work  
%% ====================================================================
\section{Background and Related Work}
\label{sec:background}

\subsection{Vehicle Routing with Synchronization Constraints}
Home healthcare routing and scheduling has emerged as a major optimization challenge due to population aging, the increasing prevalence of chronic diseases, and the growing adoption of decentralized healthcare services \citep{masmoudi2024home}. In such environments, caregivers must travel between geographically dispersed patients while respecting operational constraints including service time windows, staff qualifications, workload balancing, and route duration limitations.

Many home healthcare activities additionally require coordinated interventions involving multiple caregivers. Examples include treatments requiring complementary medical skills, patient transfers, rehabilitation sessions, and procedures involving specialized equipment. These situations introduce synchronization constraints, creating temporal dependencies between visits that must be performed simultaneously or according to predefined precedence relationships.

The Vehicle Routing Problem with Time Windows and Synchronization constraints (VRPTW-S) models these operational requirements by combining routing, scheduling, and synchronization decisions within a unified optimization framework. Compared with the classical Vehicle Routing Problem with Time Windows (VRPTW), synchronization constraints significantly increase problem complexity by coupling routes that would otherwise be independent. Consequently, modifications performed on a given route may propagate temporal infeasibilities to several other routes simultaneously, making feasibility management considerably more difficult.

The VRPTW-S belongs to the class of NP-hard combinatorial optimization problems \citep{lenstra1981complexity}. Exact optimization methods rapidly become computationally expensive as instance size and synchronization density increase, motivating the development of advanced heuristic and metaheuristic approaches capable of producing high-quality solutions within practical computational times.

\subsection{Metaheuristic Approaches for VRPTW-S}
Early research on the VRPTW-S primarily focused on exact optimization techniques. Among the pioneering contributions, \citet{bredstrom2008combined} proposed a branch-and-price framework for synchronized home healthcare routing and introduced benchmark instances that subsequently became standard references in the literature. While exact methods can provide optimal solutions for small and medium-sized instances, their scalability remains limited due to the strong temporal dependencies induced by synchronization constraints.

As problem sizes increased and synchronization requirements became more prevalent in realistic healthcare applications, research progressively shifted toward heuristic and metaheuristic approaches capable of balancing solution quality and computational effort. Among the earliest successful methods, \citet{afifi2016heuristic} developed a synchronization-aware Simulated Annealing framework incorporating specialized neighborhood operators and adaptive penalty mechanisms. Their results highlighted the importance of explicitly considering synchronization constraints during local search exploration.

Subsequent work further improved solution quality through adaptive search mechanisms. \citet{liu2019adaptive} proposed an Adaptive Large Neighborhood Search (ALNS) algorithm combining destroy-and-repair operators specifically designed for synchronized routing. By dynamically adjusting operator scores according to their historical performance, the method achieved several improvements over previously reported benchmark results.

Among neighborhood-based metaheuristics, Variable Neighborhood Search (VNS) and its extensions have demonstrated particularly strong performance. The Generalized Variable Neighborhood Search (GVNS) introduced by \citet{masmoudi2025generalized} currently represents one of the most competitive approaches reported for the VRPTW-S. Their framework combines shaking procedures, Variable Neighborhood Descent (VND), and adaptive penalty adjustment mechanisms to efficiently explore both feasible and infeasible regions of the search space. Extensive computational experiments reported numerous best-known solutions across standard benchmark sets, establishing GVNS as a strong reference method for synchronization-constrained routing optimization.

Despite the significant progress achieved by modern metaheuristics, most existing approaches continue to rely on search-control decisions that remain fixed throughout optimization. Neighborhood exploration orders, local search strategies, acceptance mechanisms, and algorithmic parameters are generally determined offline and remain unchanged regardless of the current search state. Such static decisions may limit adaptability to the highly dynamic search landscape induced by synchronization constraints.

\subsection{Learning-Guided Optimization for Routing Problems}
Recent advances in machine learning have stimulated growing interest in learning-guided combinatorial optimization. Rather than relying exclusively on handcrafted search policies, these approaches seek to learn adaptive decision mechanisms capable of exploiting information gathered during the optimization process. Comprehensive surveys of this emerging field are provided by \citet{karimi2022machine}, \citet{mazyavkina2021reinforcement}, and \citet{bengio2021machine}.

Early studies primarily focused on learning complete solution construction policies. \citet{bello2017neural} demonstrated that reinforcement learning could be used to generate solutions for the Traveling Salesman Problem without explicit supervision. This idea was subsequently extended to vehicle routing problems by \citet{nazari2018reinforcement}, who proposed a reinforcement learning framework for the Capacitated Vehicle Routing Problem, and by \citet{kool2019attention}, who introduced attention-based architectures capable of producing competitive routing solutions across several combinatorial optimization tasks.

More recently, learning has increasingly been used not to replace optimization algorithms, but to guide their search behavior. Existing approaches can broadly be categorized according to the level at which learning influences the optimization process. Some methods focus on adaptive operator selection, others target parameter control, move evaluation, acceptance decisions, or search strategy adaptation. This direction is particularly attractive for complex constrained routing problems because it preserves the strengths of established metaheuristics while introducing adaptive decision-making capabilities.

Several studies have demonstrated the effectiveness of learning-guided search control in routing optimization. For example, \citet{xu2023learning} proposed a reinforcement learning framework for adaptive neighborhood selection within a Variable Neighborhood Search procedure, showing that search performance can be improved through context-aware operator selection. Reinforcement learning has also been investigated for dynamic parameter adaptation, where policy-gradient methods and Proximal Policy Optimization (PPO) have shown promising results for controlling search behavior in complex optimization environments \citep{schulman2017proximal, wang2024csppo}. Similarly, supervised learning models have been employed to estimate move quality, predict operator effectiveness, and support adaptive search decisions during optimization.

Although these studies demonstrate the potential of learning-guided optimization, the vast majority have focused on classical routing variants that do not include synchronization constraints. Research investigating learning-based search adaptation for synchronization-constrained vehicle routing remains extremely limited, and the interaction between learning mechanisms and synchronization-aware metaheuristics remains largely unexplored.

\subsection{Research Gap and Contributions}
The literature highlights two complementary observations. First, state-of-the-art approaches for the VRPTW-S are predominantly based on advanced metaheuristics, with GVNS currently representing one of the most effective frameworks reported in the literature \citep{masmoudi2025generalized}. Second, recent studies in combinatorial optimization have demonstrated that learning-guided search control can successfully improve decision making across several routing variants.

However, these two research directions have remained largely disconnected. Existing VRPTW-S approaches continue to rely on static search-control decisions at multiple levels of the optimization process. Neighborhood operators are explored according to predefined rules, local search strategies remain fixed throughout execution, and algorithm parameters are typically tuned offline. Whether adaptive learning mechanisms can effectively replace these static decisions in synchronization-constrained routing remains largely unanswered.

To the best of our knowledge, learning-based search adaptation has received very limited attention in the VRPTW-S literature, and no previous work has investigated the combined integration of reinforcement learning and supervised learning mechanisms within a synchronization-aware GVNS framework.

This work addresses this gap through a two-stage contribution. First, we introduce GVNS+, an enhanced implementation of the state-of-the-art GVNS framework incorporating additional neighborhood operators and a new family of large-scale synchronized benchmark instances extending existing literature datasets up to 400 customers. Since existing synchronized benchmark sets remain relatively limited in scale, these larger instances were generated to enable a more comprehensive assessment of scalability, robustness, and learning generalization. The enhanced implementation is validated against the benchmark results reported by \citet{masmoudi2025generalized}, achieving an average gap of 1.14\% despite being entirely reimplemented in Python.

Second, we investigate the impact of learning-guided search adaptation on top of this enhanced baseline. Three complementary adaptive mechanisms are studied: reinforcement learning for neighborhood operator selection, supervised learning for dynamic First Improvement versus Best Improvement strategy selection, and Proximal Policy Optimization for online parameter control. Beyond evaluating the effectiveness of individual learning techniques, the study examines how the benefits of learning vary across different decision layers of the metaheuristic search process, providing insights into where adaptive intelligence can most effectively complement synchronization-aware optimization.

%% ==================================================================== 
%% 3. Problem definition
%% ====================================================================
\section{Problem Definition} 
\label{sec:problem}

\subsection{VRPTW-S formulation}
The Vehicle Routing Problem with Time Windows and Synchronization constraints (VRPTW-S) is defined on a directed graph $G=(V,E)$, where $V={0,1,\dots,n}$ denotes the set of vertices and $E$ the set of edges. Vertex $0$ represents the depot, while each customer $i \in V\setminus{0}$ requires a service of duration $s_i$ to be performed within a time window $[e_i,l_i]$.

A fleet of vehicles is available to serve all customers. Traveling from customer $i$ to customer $j$ incurs a travel time $t_{ij}$ and an associated routing cost $c_{ij}$. Each customer must be visited exactly once, and every route must start and end at the depot.

In addition to classical routing and scheduling constraints, a subset of customers requires synchronized service. Let $P^{sync}$ denote the set of synchronized visit pairs. For each pair $(i,j)\in P^{sync}$, the service start times must satisfy
\[ T_i = T_j \]
where $T_i$ and $T_j$ denote the service start times at customers $i$ and $j$, respectively.

The objective is to determine a set of feasible routes minimizing the total routing cost while satisfying vehicle capacity, route duration, time-window, and synchronization constraints. Synchronization requirements create strong temporal dependencies between routes and significantly increase the complexity of feasibility management during optimization.

\subsection{Benchmark instances}
The experimental study relies on two complementary benchmark families.

The first consists of the synchronized Solomon-derived benchmark instances commonly used in the VRPTW-S literature. These instances originate from the benchmark generation procedures introduced by \citet{bredstrom2008combined} and subsequently extended by several authors. They provide the standard basis for evaluating and comparing synchronization-aware optimization methods.

The second benchmark family was developed specifically for this work to evaluate scalability beyond the instance sizes commonly considered in the literature. Starting from classical Solomon instances, additional synchronized visits were generated following the synchronization generation principles used in existing benchmark sets. The resulting dataset includes instances containing up to 400 customers and substantially higher synchronization densities than those found in traditional benchmarks.

While the literature benchmarks are used to validate the proposed GVNS+ implementation against previously reported results, all learning-guided experiments are conducted on the newly generated benchmark family. This experimental protocol allows the adaptive mechanisms to be evaluated independently from existing benchmark-specific tuning effects while providing a more challenging assessment of scalability and robustness.


%% ====================================================================
%% 4. GVNS+ baseline framework
%% ====================================================================
\section{GVNS+: Enhanced baseline metaheuristic framework}
\label{sec:baseline}

\subsection{Original GVNS framework}
The foundation of the proposed framework is the Generalized Variable Neighborhood Search (GVNS) algorithm recently introduced by \citet{masmoudi2025generalized}, which currently represents one of the most competitive metaheuristic approaches for the Vehicle Routing Problem with Time Windows and Synchronization constraints (VRPTW-S).

The method combines diversification and intensification mechanisms through three main components. First, a shaking phase applies progressively stronger perturbations to escape local optima and diversify the search trajectory. Second, a Variable Neighborhood Descent (VND) procedure explores a sequence of neighborhood structures in order to intensify the search around promising solutions. Finally, an adaptive acceptance mechanism allows the algorithm to temporarily explore infeasible regions while dynamically regulating the penalties associated with synchronization and temporal violations.

A major strength of the GVNS lies in its ability to efficiently balance exploration and exploitation in highly constrained synchronized routing problems. By combining multiple neighborhood structures with adaptive penalty management, the method achieved state-of-the-art performance on the benchmark instances introduced by \citet{bredstrom2008combined} and later extended by \citet{liu2019adaptive}, establishing numerous best-known solutions \citep{masmoudi2025generalized}.

For these reasons, GVNS was selected as the optimization backbone of the present study. Rather than replacing the underlying metaheuristic, our objective is to strengthen its search capabilities and subsequently investigate how learning-based mechanisms can further improve its adaptive behavior.

\subsection{Added neighborhood operators}
Although the original GVNS exhibits strong performance, the present work required a robust and scalable implementation capable of supporting subsequent learning-based extensions. During the initial Python reimplementation of the original framework, a noticeable performance gap remained with respect to the reference C++ implementation reported in the literature. To strengthen the optimization backbone before introducing adaptive mechanisms, the neighborhood portfolio was extended through the integration of four additional operators.

The original framework relies on four classical neighborhoods based on relocation and exchange moves performed either within a route or between routes. To complement these local modifications, four larger-scale operators were incorporated: Cross-Exchange, $k$-Relocate, 2-Opt, and 2-Opt*.

Cross-Exchange simultaneously swaps route segments between two different routes, enabling substantial route restructuring. The $k$-Relocate operator transfers consecutive customer sequences rather than individual visits, allowing larger route modifications in a single move. The 2-Opt operator reverses route segments within a route to eliminate inefficient crossings and improve route geometry, while 2-Opt* exchanges route tails between two routes and has proven particularly effective in vehicle routing applications.

These operators were selected because of their demonstrated effectiveness in routing literature and their complementary exploration behavior \citep{liu2019adaptive}. Whereas the original neighborhoods mainly perform localized modifications, the additional operators enable larger structural transformations capable of escaping deeper local optima. The resulting framework, denoted GVNS+, preserves the overall architecture of the original algorithm while providing a richer neighborhood portfolio and enhanced search diversification capabilities.

\subsection{Large-scale benchmark generation}
Existing VRPTW-S benchmark sets provide valuable evaluation references but remain relatively limited in scale for studying adaptive optimization mechanisms. The benchmark proposed by \citet{bredstrom2008combined} contains instances derived from Solomon problems with up to 80 customers, while the larger dataset introduced by \citet{liu2019adaptive} extends problem sizes to 200 customers.

Such benchmarks are sufficient for evaluating conventional metaheuristics but provide limited opportunities for assessing scalability and learning generalization. Since the objective of this work is to investigate adaptive search mechanisms, a broader set of optimization scenarios was required.

To address this limitation, a new family of synchronized VRPTW-S instances was generated from the classical Solomon benchmark. The original customer coordinates, demands, service durations, vehicle capacities, and time-window characteristics were preserved in order to maintain the structural properties of the underlying VRPTW instances. Synchronization constraints were then introduced by creating pairs of customers requiring simultaneous service.

Four instance sizes were considered: 50, 100, 200, and 400 customers. The first two correspond to commonly used benchmark sizes, whereas the latter two significantly extend the scale of synchronized routing problems considered in the literature. For each size category, 56 synchronized instances were generated, resulting in a total collection of 224 benchmark instances.

To maintain a comparable workload across instance sizes, fleet sizes were scaled proportionally with the number of customers. Instances containing 50 and 100 customers were assigned 25 vehicles, 200-customer instances used 50 vehicles, and 400-customer instances employed 100 vehicles. Similarly, synchronization density was preserved by introducing 5, 10, 20, and 40 synchronized customers for the four size categories respectively.

Following the construction principles adopted in existing synchronized benchmarks, synchronization constraints were preferentially assigned to customers possessing relatively wide time windows. This strategy increases the likelihood of generating feasible synchronized services while preserving the routing complexity induced by temporal coordination.

The resulting benchmark extends existing VRPTW-S datasets both in scale and complexity and serves as the primary experimental environment for evaluating the adaptive mechanisms proposed in this paper.

\subsection{Learning Infrastructure and Search-State Representation}
The adaptive mechanisms introduced later in this work require the ability to observe and characterize the current optimization context. Unlike conventional machine learning applications, no training dataset exists a priori for the VRPTW-S. The information required by the learning models must therefore be generated directly from the execution of the optimization algorithm itself. To address this challenge, a unified learning infrastructure was integrated into GVNS+, enabling the continuous observation of the search process and the construction of datasets suitable for the different adaptive components investigated in this study.

Figure~\ref{fig:data_collection_pipeline} summarizes the overall data collection process. Starting from repeated executions of GVNS+ on the training benchmark instances, the framework successively performs search-state extraction, trajectory logging, offline dataset construction, and task-specific preprocessing before producing the final datasets used by the learning models.

\begin{figure}[ht]
\centering
\includegraphics[width=\linewidth]
{../figures/data_collection_pipeline.pdf}

\caption{Data collection pipeline for learning-based adaptation}
\label{fig:data_collection_pipeline}
\end{figure}

Since complete routing solutions are high-dimensional and computationally expensive to process directly, the optimization state is represented through a compact set of descriptors summarizing the current search situation. The proposed representation combines four complementary categories of features.

Instance-level descriptors characterize the solved problem through information such as instance size, fleet size, synchronization density, planning horizon, and time-window characteristics. These features remain constant during a run and provide contextual information regarding the difficulty and structure of the optimization instance.

Dynamic search descriptors capture the evolution of the optimization process over time. They include indicators related to iteration progress, search stagnation, recent improvement rates, and the evolution of solution quality over a rolling observation window. These features provide information about whether the search is actively progressing or entering a stagnation phase.

Solution-structure descriptors summarize the current routing solution through lightweight indicators reflecting route organization and feasibility pressure. Examples include route-balance measures, waiting-time statistics, lateness indicators, and penalty-related information. Such descriptors provide an approximate characterization of solution quality while remaining computationally inexpensive.

Finally, neighborhood-context descriptors characterize the local search environment by recording information related to the active neighborhood operator and its recent behavior. These descriptors allow the adaptive models to distinguish decisions occurring under different exploration conditions.

To generate learning data, GVNS+ was executed repeatedly on the training benchmark instances. During each run, the search-state descriptors were continuously recorded together with the corresponding optimization decisions and their outcomes. Because the metaheuristic contains stochastic components, repeated executions of the same instance naturally generate diverse optimization trajectories, exposing the learning models to a wide range of search situations, including productive improvement phases, neighborhood transitions, and prolonged stagnation periods.

The collected logs were subsequently processed offline to construct task-specific learning datasets. Depending on the considered adaptive mechanism, additional labels, rewards, or decision targets were derived from the recorded optimization outcomes. Basic preprocessing operations, including duplicate filtering, removal of non-informative variables, and feature-correlation analysis, were applied to reduce redundancy while preserving the most relevant information for each learning task.

The resulting infrastructure provides a common data foundation for all adaptive components developed in this work. Although the learning objectives differ across reinforcement learning, supervised learning, and policy-learning approaches, all models operate from the same unified representation of the search process, ensuring consistency while maintaining a compact and computationally efficient description of the optimization context.
\subsection{Validation against literature benchmarks}
Before introducing learning-based adaptations, the performance of GVNS+ was evaluated on the benchmark instances commonly used in the VRPTW-S literature. The objective was to verify that the proposed enhancements preserve the effectiveness of the original method while reducing the performance gap introduced during the Python reimplementation.

The validation experiments were conducted on the benchmark instances of \citet{bredstrom2008combined} and \citet{liu2019adaptive}, using the same evaluation protocol adopted by \citet{masmoudi2025generalized}. The obtained solutions were compared against the best-known results reported in the literature.

The integration of the additional neighborhood operators substantially improved solution quality across all benchmark categories. Whereas the initial reimplementation exhibited an average performance gap of approximately 10\% relative to the reference implementation, the enhanced GVNS+ reduced this gap to 1.14\% on average. Furthermore, the method remained highly competitive with respect to published state-of-the-art results while preserving reasonable computational times.

These results confirm that the proposed enhancements successfully strengthen the optimization backbone without altering the fundamental search principles of the original GVNS. Consequently, GVNS+ provides a reliable and competitive baseline from which the impact of the subsequent learning-guided adaptations can be assessed.


%% ====================================================================
%% 5. Proposed learning-guided adaptive framework
%% ====================================================================
\section{Learning-Guided Search Adaptation} 
\label{sec:framework}

%% Adaptive neighborhood ordering
\subsection{RL-based operator selection} 
The performance of Variable Neighborhood Descent depends not only on the quality of the available neighborhood operators but also on the order in which they are explored. In the original GVNS and in the enhanced GVNS+ framework, neighborhoods are evaluated according to a predefined sequence composed of insertion, swap, relocation, exchange, and route-improvement operators. Although such a deterministic strategy provides stable performance, it assumes that the same neighborhood ordering remains appropriate throughout the entire optimization process.

In practice, the effectiveness of a neighborhood strongly depends on the current search context. Large structural moves may be beneficial immediately after diversification phases, whereas smaller operators often become more effective near local optima. Similarly, the most promising neighborhood may vary according to feasibility status, search stagnation, instance characteristics, and the recent evolution of solution quality. Consequently, a fixed neighborhood ordering may fail to exploit the search dynamics efficiently.

To address this limitation, neighborhood selection is formulated as a sequential decision-making problem solved through reinforcement learning. Instead of exploring neighborhoods according to a predefined order, the algorithm dynamically selects the operator that appears most promising given the current optimization state.

\subsubsection{MDP Formulation}
The adaptive operator-selection problem is modeled as a Markov Decision Process (MDP) embedded within the BVND intensification phase.

At each decision step, the learning agent observes a compact representation of the current optimization context and selects one neighborhood operator from the GVNS+ portfolio. The selected operator is applied to the current solution, producing a new search state and a reward that reflects the quality of the obtained improvement.

The state representation follows the unified search-state formulation introduced in Section~\ref{sec:baseline}. Rather than encoding complete routing solutions, which would be computationally prohibitive, the search process is described through a set of aggregated descriptors characterizing four complementary aspects of the optimization context:

\begin{itemize}
\item instance-level characteristics,
\item current solution quality and feasibility pressure,
\item search dynamics and stagnation indicators,
\item neighborhood effectiveness statistics.
\end{itemize}

This representation provides sufficient information for decision making while remaining computationally lightweight.

The action space corresponds directly to the eight neighborhood operators available in GVNS+:

\begin{equation}
A=\{N_1,N_2,\ldots,N_8\}.
\end{equation}

Each action represents the selection of one neighborhood operator for the next local-search step.

The reward is based on the improvement achieved in the augmented objective function. Let $F_t$ denote the objective value before applying the selected operator and $F_{t+1}$ the value obtained afterward. The normalized improvement is defined as

\begin{equation}
r_t=\frac{F_t-F_{t+1}}{1+F_t}.
\end{equation}

Additional reward shaping terms are introduced to encourage feasibility improvements and discourage ineffective neighborhood selections. Detailed reward definitions for each learning model are provided in Appendix~B.

Following the application of the selected operator, the optimization context is updated and becomes the next state observed by the agent. The learning process continues until the termination criterion of the outer GVNS+ search is reached.

\subsubsection{Integration within GVNS+}
Figure~\ref{fig:rl_gvns_architecture} illustrates the integration of reinforcement learning within the GVNS+ framework.

\begin{figure}[ht]
\centering
\includegraphics[width=\linewidth]
{../figures/rl_gvns_architecture.pdf}

\caption{RL-driven GVNS+ architecture}
\label{fig:rl_gvns_architecture}
\end{figure}

At each GVNS+ iteration, the shaking phase first generates a perturbed solution intended to escape the current local optimum. The search-state representation is then extracted and provided to the RL agent. Based on this information, the agent selects one neighborhood operator to be executed during the subsequent intensification step.

The selected neighborhood is applied through a single-neighborhood BVND iteration, generating a candidate solution. The environment then evaluates the resulting improvement, computes the reward, and updates the state representation. This transition is used by the learning algorithm to refine its decision policy. Finally, the candidate solution is processed by the GVNS+ acceptance mechanism before the search proceeds to the next iteration.

This integration preserves the overall structure of GVNS+ while replacing the static neighborhood-ordering mechanism with an adaptive decision process capable of reacting to the current optimization context.

\subsubsection{Tabular Q-Learning Agent}
The first adaptive operator-selection strategy is based on tabular Q-learning \cite{sutton2018reinforcement}. The objective is to learn the expected long-term utility of each neighborhood operator under different search conditions.

Because tabular reinforcement learning requires a finite state space, the continuous search descriptors are transformed into a compact abstract representation. The abstraction combines feasibility indicators, stagnation information, neighborhood-performance statistics, and instance-size categories. The resulting state space remains sufficiently small to enable efficient learning while preserving the main characteristics of the optimization context.

\input{../tables/qlearning_state.tex}

The agent maintains a Q-table storing an estimate of the expected cumulative reward associated with each state-action pair. After observing a transition $(s,a,r,s')$, the table is updated according to the classical Bellman update rule:

\begin{equation}
Q(s,a)\leftarrow Q(s,a)+
\alpha
\left[
r+\gamma \max_{a'}Q(s',a')
-Q(s,a)
\right].
\end{equation}

Action selection is performed using an Upper Confidence Bound (UCB) strategy, which balances exploration and exploitation by considering both the estimated value of an action and the number of times it has been previously selected. Compared with standard $\epsilon$-greedy exploration, UCB promotes a more structured exploration process and is particularly suitable for adaptive metaheuristics where neighborhood effectiveness evolves throughout the search.

\subsubsection{Deep Q-Network Agent}
Although tabular Q-learning provides an interpretable adaptive mechanism, its reliance on state discretization limits its ability to capture fine-grained variations in the optimization process. To overcome this limitation, a Deep Q-Network (DQN) \cite{mnih2015human} is also investigated.

Unlike tabular Q-learning, DQN approximates the action-value function through a neural network operating directly on the continuous search-state representation. This eliminates the need for manual state discretization and enables the model to learn more complex relationships between optimization states and neighborhood effectiveness.

The adopted network architecture is shown in Figure~\ref{fig:dqn_architecture}. The model consists of a multilayer perceptron with two hidden layers of 128 neurons each and ReLU activation functions. The output layer contains one neuron for each neighborhood operator and produces the estimated action values $Q(s,a)$.

\begin{figure}[ht]
\centering
\includegraphics[width=\linewidth]
{../figures/dqn_architecture.pdf}

\caption{DQN architecture}
\label{fig:dqn_architecture}
\end{figure}

To improve learning stability across instances of different sizes, all state features are normalized online using running estimates of feature means and variances before being processed by the network. This prevents large-magnitude descriptors from dominating the learning process and improves generalization across heterogeneous routing instances.

Action selection follows an $\epsilon$-greedy exploration strategy. With probability $\epsilon$, a neighborhood is selected randomly; otherwise, the neighborhood associated with the largest predicted action value is chosen.

Experience replay and target-network mechanisms are employed during training to stabilize value estimation and reduce temporal correlations between consecutive observations. The network parameters are optimized by minimizing the temporal-difference error over mini-batches sampled from the replay buffer.

% Adaptive FI/BI selection
\subsection{RF-based strategy selection} 
\label{sec:rf_selector}

The efficiency of the B-VND intensification phase depends not only on the selected neighborhood operator but also on the exploration strategy used within that neighborhood. In the original GVNS framework, local search is performed using a First Improvement (FI) strategy, where candidate moves are evaluated sequentially and the search terminates as soon as an improving move is identified. This approach is computationally efficient and allows rapid progression through large neighborhoods.

An alternative strategy is Best Improvement (BI), which evaluates all admissible moves before selecting the most profitable one. Although BI generally provides stronger intensification, it incurs substantially higher computational costs, particularly for expensive neighborhoods such as cross-exchange and 2-opt* on large synchronized instances.

Preliminary analyses conducted on GVNS+ indicated that neither strategy consistently dominates the other. FI is generally advantageous during early search phases, where improving moves are abundant and rapid progress is desirable. Conversely, BI becomes more beneficial near local optima, where improvements are rarer and a more exhaustive exploration can uncover higher-quality moves. The effectiveness of each strategy therefore depends on the current optimization context, including the search stage, neighborhood type, stagnation level, and solution structure.

Rather than selecting a single exploration strategy for the entire optimization process, the problem is reformulated as a contextual decision task in which the algorithm dynamically chooses between FI and BI according to the current search state.

%% ======== Two-head RF selector architecture =========
\subsubsection{Two-head Random Forest strategy selector}
The proposed selector is based on a two-head Random Forest (RF) architecture. Both heads operate on the unified search-state representation introduced in Section~\ref{sec:learning_infra}, which combines instance descriptors, dynamic search indicators, solution characteristics, and neighborhood-context information.

The first head is formulated as a binary classification task estimating the probability that a given strategy will produce an improving move:

\begin{equation}
P(\text{success}\mid s,\text{strategy}),
\label{eq:rf_prob}
\end{equation}

where a successful move corresponds to a strictly positive improvement of the augmented objective function.

The second head is formulated as a regression task estimating the expected normalized improvement:

\begin{equation}
\Delta_{\text{rel}}
=
\frac{F_{\text{before}}-F_{\text{after}}}
{1+F_{\text{before}}},
\label{eq:rf_gain}
\end{equation}

where $F$ denotes the augmented objective function.

This decomposition enables the model to distinguish between the reliability of a strategy and the magnitude of the expected improvement. A strategy may frequently produce small gains or occasionally generate larger improvements, and modeling these aspects separately improves both interpretability and robustness.

%% ======== Training procedure =========
\subsubsection{Training procedure}
Training data are generated using the learning infrastructure presented previously. For each training instance, GVNS+ is executed twice: once using FI exclusively and once using BI exclusively. During these executions, all local-search decisions are recorded together with the corresponding search-state descriptors and optimization outcomes.

Each recorded observation is transformed into a supervised-learning sample by augmenting the search-state representation with an indicator specifying the evaluated strategy. The resulting dataset is then used to train the classifier and regressor independently.

The classifier learns to predict the probability of obtaining an improving move, whereas the regressor estimates the expected normalized improvement. Both models use identical Random Forest configurations, whose hyperparameters are reported in Appendix~B.

%% ======== Online decision mechanism =========
\subsubsection{Online decision mechanism}
During optimization, the selector evaluates the current search state under two hypothetical scenarios corresponding to FI and BI.

Given a feature vector $\mathbf{x}$ describing the current optimization context, two augmented representations are generated:

\begin{equation}
\mathbf{x}_{FI}
=
[\mathbf{x},0],
\label{eq:xfi}
\end{equation}

\begin{equation}
\mathbf{x}_{BI}
=
[\mathbf{x},1].
\label{eq:xbi}
\end{equation}

The classifier and regressor evaluate both representations and produce four predictions:

\[
P_{FI},\quad G_{FI},\quad
P_{BI},\quad G_{BI},
\]

where $P$ denotes the predicted success probability and $G$ the expected normalized gain.

The final score associated with each strategy is computed as

\begin{equation}
S_{FI}
=
P_{FI}\max(G_{FI},0),
\label{eq:sfi}
\end{equation}

\begin{equation}
S_{BI}
=
P_{BI}\max(G_{BI},0).
\label{eq:sbi}
\end{equation}

Negative predicted gains are clipped to zero in order to avoid rewarding deteriorating moves. The strategy associated with the largest score is selected for neighborhood exploration.

This multiplicative scoring rule naturally balances reliability and expected improvement magnitude. Strategies that appear promising but unreliable receive lower scores, while highly reliable strategies producing only negligible gains are also penalized.

Figure~\ref{fig:rf_internal_layout} and Figure~\ref{fig:rf_inference} summarize the complete online inference procedure.

\begin{figure}[ht]
\centering
\includegraphics[width=\linewidth]
{../figures/rf_internal_layout.pdf}

\caption{RF internal layout}
\label{fig:rf_internal_layout}
\end{figure}

\begin{figure}[ht]
\centering
\includegraphics[width=\linewidth]
{../figures/rf_inference.pdf}

\caption{RF inference procedure}
\label{fig:rf_inference}
\end{figure}

%% ========= Integration into GVNS+ =========
\subsubsection{Integration into GVNS+}
The trained RF models are loaded once at the beginning of each optimization run. Before exploring a neighborhood during the B-VND procedure, the current search-state features are extracted and evaluated by the selector. The resulting decision determines whether the neighborhood is explored using FI or BI.

Algorithm~\ref{alg:rf_selector} summarizes the complete integration procedure.

\begin{algorithm}[ht]
\caption{RF-Based Adaptive Strategy Selection}
\label{alg:rf_selector}
\begin{algorithmic}[1]

\Require Current solution $S$, neighborhood $N_k$
\Ensure Candidate solution $S'$

\State Extract search-state features $\mathbf{x}$

\State Construct
$\mathbf{x}_{FI}=[\mathbf{x},0]$

\State Construct
$\mathbf{x}_{BI}=[\mathbf{x},1]$

\State Predict
$P_{FI}$ and $G_{FI}$

\State Predict
$P_{BI}$ and $G_{BI}$

\State Compute
$S_{FI}=P_{FI}\max(G_{FI},0)$

\State Compute
$S_{BI}=P_{BI}\max(G_{BI},0)$

\If{$S_{BI} > S_{FI}$}
    \State Apply neighborhood $N_k$ using BI
\Else
    \State Apply neighborhood $N_k$ using FI
\EndIf

\State Return resulting solution

\end{algorithmic}
\end{algorithm}

The proposed mechanism introduces an adaptive intensification layer within GVNS+, allowing the local-search strategy to evolve according to the current optimization context rather than remaining fixed throughout the search process.


% PPO-based parameter control
\subsection{PPO-based parameter control}
\label{sec:ppo}

Although GVNS+ provides a stronger optimization backbone than the original GVNS, its overall behaviour still depends on a set of manually calibrated hyperparameters controlling diversification, intensification, feasibility management, and stagnation recovery. In the original framework, these parameters are fixed before the search through offline calibration \cite{masmoudi2025generalized}. However, the most appropriate parameter configuration may vary significantly throughout the optimization process. Early search stages often benefit from stronger diversification, whereas later stages generally require more intensive exploitation around promising regions of the search space.

To address this limitation, we introduce a reinforcement learning controller based on Proximal Policy Optimization (PPO) \cite{schulman2017proximal}. Unlike the operator-selection agents presented previously, which act locally inside the B-VND procedure, the PPO controller operates at the global level of the metaheuristic by dynamically adapting the main search-control parameters during optimization.

%% ========= MDP formulation =========
\subsubsection{MDP Formulation}
The adaptive hyperparameter control problem is modeled as a Markov Decision Process (MDP). At each decision point, the controller observes the current optimization context, selects a search-control action, and receives a reward reflecting the impact of the resulting parameter configuration on search performance.

%% state represententaion
\paragraph{State representation}
The controller relies on the unified search-state representation introduced in Section~\ref{sec:learninginfra}. The state combines instance-level descriptors, dynamic search indicators, solution-structure characteristics, and neighborhood-context information.

In addition, the current values of the controlled hyperparameters are included in the observation vector. This allows the policy to reason jointly about the optimization state and the current search configuration. The resulting state can be expressed as

\[
s_t=
[x_t^{inst},
x_t^{search},
x_t^{sol},
x_t^{neigh},
x_t^{ctrl}]
\]

where the last component corresponds to the current parameter values.

%% action space
\paragraph{Action space}
Rather than directly predicting numerical parameter values, the controller selects among a set of high-level search behaviours. Each action modifies several parameters simultaneously and corresponds to a meaningful optimization strategy.

The controller regulates four key components of GVNS+:

\begin{itemize}
\item shaking intensity,
\item acceptance threshold,
\item penalty adaptation scale,
\item stagnation threshold.
\end{itemize}

To avoid unstable parameter oscillations, updates are smoothed using an inertia mechanism:

\[
\theta_{new}
=
0.7\theta_{current}
+
0.3\theta_{target}.
\]

The resulting values are subsequently clipped to predefined safety ranges.

\input{../tables/ppo_parameter_ranges}

The complete action space is summarized in Table~\ref{tab:ppo_actions}. Each action corresponds to a distinct search behaviour ranging from aggressive diversification to feasibility reinforcement and stagnation recovery.

\input{../tables/ppo_actions}

%% reward function
\paragraph{Reward function}

The reward is designed to measure the quality of the selected parameter configuration while maintaining numerical stability during policy optimization.

Let $F_{before}$ and $F_{after}$ denote the penalized objective values before and after the corresponding local-search step. The relative improvement is computed as

\[
r_{imp}
=
\frac{F_{before}-F_{after}}
{\max(1,F_{before})}.
\]

To reduce sensitivity to extreme values, the improvement is transformed using

\[
r_{base}
=
\tanh(3r_{imp}).
\]

Additional shaping terms reward improving moves and newly discovered best solutions while slightly penalizing non-improving decisions. The final reward is clipped to the interval $[-1,1]$.

This formulation provides stable learning signals while preserving sensitivity to meaningful improvements in solution quality.

%% ========= Architecture =========
\subsubsection{PPO Controller Architecture}

The controller is implemented as a compact actor--critic neural network. Since the state representation is already engineered and normalized, a lightweight architecture is sufficient and avoids unnecessary model complexity.

The network consists of two fully connected hidden layers of 128 neurons with ReLU activations. These layers produce a shared latent representation

\[
h=f_{\theta}(s).
\]

Two output heads are then derived from this representation.

The actor head produces a probability distribution over the eight control actions:

\[
\pi(a|s)=
\mathrm{softmax}(W_{\pi}h).
\]

The critic head estimates the state value

\[
V(s)=W_v h.
\]

The actor is responsible for selecting parameter-control actions, whereas the critic estimates the expected future return and serves as a baseline for variance reduction during training.

\begin{figure}[ht]
\centering
\includegraphics[width=\linewidth]
{../figures/ppo_architecture.pdf}

\caption{PPO architecture}
\label{fig:ppo_architecture}
\end{figure}

%% ========= Training =========
\subsubsection{Offline Policy Training}

Training data are obtained from the learning infrastructure presented previously. During repeated executions of GVNS+, optimization trajectories are recorded as sequences of transitions

\[
(s_t,a_t,r_t,s_{t+1}).
\]

The PPO controller is subsequently trained offline using these trajectories.

For each trajectory, discounted returns are computed according to

\[
G_t=r_t+\gamma G_{t+1},
\]

with $\gamma=0.99$. Advantages are then estimated as

\[
A_t=G_t-V(s_t).
\]

The policy is optimized using the clipped PPO objective proposed by Schulman et al.~\cite{schulman2017proximal}, while the critic is trained through mean squared error regression on the estimated returns. An entropy regularization term is additionally included to encourage exploration during learning.

Figure~\ref{fig:ppo_training_pipeline} summarizes the complete offline training workflow.

\begin{figure}[ht]
\centering
\includegraphics[width=\linewidth]
{../figures/ppo_training_pipeline.pdf}

\caption{PPO training pipeline}
\label{fig:ppo_training_pipeline}
\end{figure}

%% ========= Integration into GVNS+ =========
\subsubsection{Integration into GVNS+}

During optimization, the controller operates as a global adaptation layer placed above the standard GVNS+ search process.

At each neighborhood evaluation, the current search-state vector is extracted and provided to the PPO policy. The selected action updates the search-control parameters, after which the corresponding local-search phase is executed using the new configuration. The resulting optimization outcome is then used to compute the reward and update the search state.

Unlike static parameter calibration, this mechanism continuously adapts the search behaviour according to the current optimization context. The controller therefore regulates the balance between diversification, intensification, feasibility restoration, and stagnation recovery throughout the entire search trajectory.

Algorithm~\ref{alg:ppo_gvns} presents the overall integration of the controller within GVNS+, while Algorithm~\ref{alg:ppo_local_search} details the PPO-guided local search procedure.

\input{../algorithms/ppo_integration}



%% =====================================================================
%% 6. RESULTS AND DISCUSSION
%% =====================================================================
\section{Experimental Setup}
\label{sec:setup}

\section{Results and Discussion}
\label{sec:results}

\subsection{GVNS+ convergence analysis}

\subsection{Operator selection}

\subsection{FI/BI selection}

\subsection{Hybrid models}

\subsection{Parameter control}

\subsection{Comparative analysis}

%% =====================================================================
%% CONCLUSION
%% =====================================================================
\section{Conclusion}
\label{sec:conclusion}

%% ----------------------------------------------------------------------
%% BIBLIOGRAPHY
%% ----------------------------------------------------------------------
\bibliographystyle{elsarticle-num-names}
\bibliography{bibliography/references}

\end{document}